\addcontentsline{toc}{chapter}{Abstract}
% A short summary of the project, the report and what the reader
% will get from reading it.


\noindent
Efficient heat production and energy management are becoming more important in today’s district heating systems, especially when operational costs, electricity prices, and energy needs keep changing. The project focuses on creating an application that improves the scheduling of heat production across several production units. The system looks at how these production units are set up, what the heat needs are and the conditions of the electricity market to figure out the most affordable way to produce energy in various operating situations.
\par
The project was done using Scrum method and by following steps to develop software incrementally. The team has used several sprints to organize tasks, manage work together and improve through meetings, reviews, and retrospectives. In addition to project management activities, the team focused on software design and system architecture. 
\par  
The document presents the entire development process, which starts with the project background, objectives and requirements and is followed by the Scrum implementation and sprint activities. It also describes how the solution is built by using UML diagrams, architectural descriptions, frontend and backend implementation details and data management strategies. In addition, teamwork, testing, refactoring, code reviews, and requirements engineering are discussed to evaluate both the technical aspects of the project and the effectiveness of team collaboration. 
\par
By looking at the different stages of development and the decisions made during the process, the reader learns how Agile methods and software engineering techniques have been used to build a working optimization system for district heating production.
